\documentclass[10pt,a4paper]{amsart}
\usepackage{amsmath,amssymb,amsthm,graphicx}
% Retained minimal preamble for the Stacks Project chapter varieties.tex.
% Verbatim excerpt of the chapter preamble: Xy-pic setup (native preamble lines 14-18), the native theorem
% declarations with the mathematical macros (native preamble lines 154-217), and the project own no-op
% definitions of the reference/slogan/history display environments (native wrapper.tex lines 1-3) that the
% counted statement uses for its slogan. Copied mechanically; nothing added, removed or redefined.
% For commutative diagrams we use Xy-pic
\usepackage[all]{xy}

% We use 2cell for 2-commutative diagrams.
\xyoption{2cell}
% Theorem environments.
%
\theoremstyle{plain}
\newtheorem{theorem}[subsection]{Theorem}
\newtheorem{proposition}[subsection]{Proposition}
\newtheorem{lemma}[subsection]{Lemma}

\theoremstyle{definition}
\newtheorem{definition}[subsection]{Definition}
\newtheorem{example}[subsection]{Example}
\newtheorem{exercise}[subsection]{Exercise}
\newtheorem{situation}[subsection]{Situation}

\theoremstyle{remark}
\newtheorem{remark}[subsection]{Remark}
\newtheorem{remarks}[subsection]{Remarks}

\numberwithin{equation}{subsection}

% Macros
%
\def\lim{\mathop{\mathrm{lim}}\nolimits}
\def\colim{\mathop{\mathrm{colim}}\nolimits}
\def\Spec{\mathop{\mathrm{Spec}}}
\def\Hom{\mathop{\mathrm{Hom}}\nolimits}
\def\Ext{\mathop{\mathrm{Ext}}\nolimits}
\def\SheafHom{\mathop{\mathcal{H}\!\mathit{om}}\nolimits}
\def\SheafExt{\mathop{\mathcal{E}\!\mathit{xt}}\nolimits}
\def\Sch{\mathit{Sch}}
\def\Mor{\mathop{\mathrm{Mor}}\nolimits}
\def\Ob{\mathop{\mathrm{Ob}}\nolimits}
\def\Sh{\mathop{\mathit{Sh}}\nolimits}
\def\NL{\mathop{N\!L}\nolimits}
\def\CH{\mathop{\mathrm{CH}}\nolimits}
\def\proetale{{pro\text{-}\acute{e}tale}}
\def\etale{{\acute{e}tale}}
\def\QCoh{\mathit{QCoh}}
\def\Ker{\mathop{\mathrm{Ker}}}
\def\Im{\mathop{\mathrm{Im}}}
\def\Coker{\mathop{\mathrm{Coker}}}
\def\Coim{\mathop{\mathrm{Coim}}}

% Boxtimes
%
\DeclareMathSymbol{\boxtimes}{\mathbin}{AMSa}{"02}

%
% Macros for moduli stacks/spaces
%
\def\QCohstack{\mathcal{QC}\!\mathit{oh}}
\def\Cohstack{\mathcal{C}\!\mathit{oh}}
\def\Spacesstack{\mathcal{S}\!\mathit{paces}}
\def\Quotfunctor{\mathrm{Quot}}
\def\Hilbfunctor{\mathrm{Hilb}}
\def\Curvesstack{\mathcal{C}\!\mathit{urves}}
\def\Polarizedstack{\mathcal{P}\!\mathit{olarized}}
\def\Complexesstack{\mathcal{C}\!\mathit{omplexes}}
% \Pic is the operator that assigns to X its picard group, usage \Pic(X)
% \Picardstack_{X/B} denotes the Picard stack of X over B
% \Picardfunctor_{X/B} denotes the Picard functor of X over B
\def\Pic{\mathop{\mathrm{Pic}}\nolimits}
\def\Picardstack{\mathcal{P}\!\mathit{ic}}
\def\Picardfunctor{\mathrm{Pic}}
\def\Deformationcategory{\mathcal{D}\!\mathit{ef}}
\newenvironment{reference}{}{}
\newenvironment{slogan}{}{}
\newenvironment{history}{}{}
\usepackage{cleveref}
\crefname{originalreference}{original reference}{original references}
\Crefname{originalreference}{Original reference}{Original references}
\makeatletter
\newlabel{limits-proposition-affine}{{05YU}{source}{}{sourcepointer-1}{}}
\newlabel{limits-proposition-affine@cref}{{[originalreference][0][]05YU}{[1][0][]source}{}{}{}}
\makeatother
\crefname{originalreference}{original reference}{original references}
\Crefname{originalreference}{Original reference}{Original references}
\makeatletter
\newlabel{limits-lemma-relative-approximation}{{09MV}{source}{}{sourcepointer-2}{}}
\newlabel{limits-lemma-relative-approximation@cref}{{[originalreference][0][]09MV}{[1][0][]source}{}{}{}}
\makeatother
\crefname{originalreference}{original reference}{original references}
\Crefname{originalreference}{Original reference}{Original references}
\makeatletter
\newlabel{limits-lemma-limit-affine}{{01Z6}{source}{}{sourcepointer-3}{}}
\newlabel{limits-lemma-limit-affine@cref}{{[originalreference][0][]01Z6}{[1][0][]source}{}{}{}}
\makeatother
\crefname{originallemma}{original lemma}{original lemmas}
\Crefname{originallemma}{Original lemma}{Original lemmas}
\makeatletter
\newlabel{lemma-local-isomorphism}{{09NH}{source}{}{sourcepointer-4}{}}
\newlabel{lemma-local-isomorphism@cref}{{[originallemma][0][]09NH}{[1][0][]source}{}{}{}}
\makeatother
\crefname{originallemma}{original lemma}{original lemmas}
\Crefname{originallemma}{Original lemma}{Original lemmas}
\makeatletter
\newlabel{lemma-characterize-open-immersion}{{09NG}{source}{}{sourcepointer-5}{}}
\newlabel{lemma-characterize-open-immersion@cref}{{[originallemma][0][]09NG}{[1][0][]source}{}{}{}}
\makeatother
\crefname{originalreference}{original reference}{original references}
\Crefname{originalreference}{Original reference}{Original references}
\makeatletter
\newlabel{properties-lemma-ample-finite-set-in-affine}{{01ZY}{source}{}{sourcepointer-6}{}}
\newlabel{properties-lemma-ample-finite-set-in-affine@cref}{{[originalreference][0][]01ZY}{[1][0][]source}{}{}{}}
\makeatother
\title{A finite set of points with matching local rings on an affine source lies in an affine open of the quasi-separated target}
\author{The Stacks Project Authors}
\date{stacks-project:varieties.tex:local-snapshot-20261011; GNU FDL 1.2 or later}
\begin{document}
\maketitle
\noindent\textit{Editorial note.} All delivered mathematical text is the author's own, copied verbatim from the pinned Stacks Project chapter file varieties.tex; no statement or proof step was rewritten, reordered or removed.\par\medskip
\noindent\textit{Editorial note.} The counted span cites six cross-chapter labels, each recorded in assembly.external\_references with its canonical tag as a navigation pointer.\par\medskip
\noindent\textit{Editorial note.} The proof cites no bibliography item inside the counted span.\par\medskip
\noindent\textit{Allowed standard prerequisites (editorial).}\begin{itemize}
\item Criterion for affineness via a limit presentation. Stacks Project, Limits of Schemes, Proposition limits-proposition-affine (tag 01Z2). Provides the affine-open criterion on which the whole proof rests.
\item Local rings and generic points under a morphism. Stacks Project, Schemes (generic points and local rings of schemes). Lets the hypotheses on T\textasciicircum{}0 -> X\textasciicircum{}0 and the local rings be used to control the points x\_j.
\end{itemize}\medskip
\ifdefined\hypertarget\hypertarget{sourcepointer-1}{}\fi
\noindent\textit{Source reference.} Original reference 05YU: sources/stacks-chow/tags.txt, source lines 7668--7668. Cross-chapter standard prerequisite cited by the author inside the counted span; it is used as a pointer and its own proof is not part of this item.\par\medskip
\ifdefined\hypertarget\hypertarget{sourcepointer-2}{}\fi
\noindent\textit{Source reference.} Original reference 09MV: sources/stacks-chow/tags.txt, source lines 12193--12193. Cross-chapter standard prerequisite cited by the author inside the counted span; it is used as a pointer and its own proof is not part of this item.\par\medskip
\ifdefined\hypertarget\hypertarget{sourcepointer-3}{}\fi
\noindent\textit{Source reference.} Original reference 01Z6: sources/stacks-chow/tags.txt, source lines 2581--2581. Cross-chapter standard prerequisite cited by the author inside the counted span; it is used as a pointer and its own proof is not part of this item.\par\medskip
\ifdefined\hypertarget\hypertarget{sourcepointer-4}{}\fi
\noindent\textit{Source reference.} Original lemma 09NH: sources/stacks-chow/tags.txt, source lines 12215--12215. Cross-chapter standard prerequisite cited by the author inside the counted span; it is used as a pointer and its own proof is not part of this item.\par\medskip
\ifdefined\hypertarget\hypertarget{sourcepointer-5}{}\fi
\noindent\textit{Source reference.} Original lemma 09NG: sources/stacks-chow/tags.txt, source lines 12214--12214. Cross-chapter standard prerequisite cited by the author inside the counted span; it is used as a pointer and its own proof is not part of this item.\par\medskip
\ifdefined\hypertarget\hypertarget{sourcepointer-6}{}\fi
\noindent\textit{Source reference.} Original reference 01ZY: sources/stacks-chow/tags.txt, source lines 2613--2613. Cross-chapter standard prerequisite cited by the author inside the counted span; it is used as a pointer and its own proof is not part of this item.\par\medskip
\begin{lemma}
\label{lemma-points-in-affine}
Let $f : T \to X$ be a morphism of schemes. Let $X^0$, resp.\ $T^0$
denote the sets of generic points of irreducible components.
Let $t_1, \ldots, t_m \in T$ be a finite set of points
with images $x_j = f(t_j)$. If
\begin{enumerate}
\item $T$ is affine,
\item $X$ is quasi-separated,
\item $X^0$ is finite
\item $f(T^0) \subset X^0$ and $f : T^0 \to X^0$ is injective, and
\item $\mathcal{O}_{X, x_j} = \mathcal{O}_{T, t_j}$,
\end{enumerate}
then there exists an affine open of $X$ containing $x_1, \ldots, x_r$.
\end{lemma}

\begin{proof}
Using Limits, Proposition \ref{limits-proposition-affine}
there is an immediate reduction to the case where $X$ and $T$ are reduced.
Details omitted.

\medskip\noindent
Assume $X$ and $T$ are reduced. We may write $T = \lim_{i \in I} T_i$
as a directed limit of schemes of finite presentation over $X$
with affine transition morphisms, see
Limits, Lemma \ref{limits-lemma-relative-approximation}.
Pick $i \in I$ such that $T_i$ is affine, see
Limits, Lemma \ref{limits-lemma-limit-affine}.
Say $T_i = \Spec(R_i)$ and $T = \Spec(R)$.
Let $R' \subset R$ be the image of $R_i \to R$.
Then $T' = \Spec(R')$ is affine, reduced, of finite type over $X$,
and $T \to T'$ dominant. For $j = 1, \ldots, r$ let $t'_j \in T'$
be the image of $t_j$. Consider the local ring maps
$$
\mathcal{O}_{X, x_j} \to
\mathcal{O}_{T', t'_j} \to
\mathcal{O}_{T, t_j}
$$
Denote $(T')^0$ the set of generic points of irreducible
components of $T'$. Let $\xi \leadsto t'_j$ be a specialization with
$\xi \in (T')^0$. As $T \to T'$ is dominant we can choose $\eta \in T^0$ 
mapping to $\xi$ (warning: a priori we do not know that $\eta$ specializes
to $t_j$). Assumption (3) applied to $\eta$ tells us that the image $\theta$
of $\xi$ in $X$ corresponds to a minimal prime of $\mathcal{O}_{X, x_j}$.
Lifting $\xi$ via the isomorphism of (5) we obtain a specialization
$\eta' \leadsto t_j$ with $\eta' \in T^0$ mapping to $\theta \leadsto x_j$.
The injectivity of (4) shows that $\eta = \eta'$. Thus
every minimal prime of $\mathcal{O}_{T', t'_j}$ lies below
a minimal prime of $\mathcal{O}_{T, t_j}$. We conclude that
$\mathcal{O}_{T', t'_j} \to \mathcal{O}_{T, t_j}$ is injective,
hence both maps above are isomorphisms.

\medskip\noindent
By Lemma \ref{lemma-local-isomorphism} there exists an open
$U \subset T'$ containing all the points $t'_j$ such that
$U \to X$ is a local isomorphism as in
Lemma \ref{lemma-characterize-open-immersion}.
By that lemma we see that $U \to X$ is an open immersion.
Finally, by
Properties, Lemma \ref{properties-lemma-ample-finite-set-in-affine}
we can find an open $W \subset U \subset T'$
containing all the $t'_j$. The image of $W$ in $X$ is the
desired affine open.
\end{proof}

\end{document}
